% lecturekit.tex
% The study PDF: Beamer theme for the lecturekit dialect. Colours are the
% nm<Name> tokens that the dual-output filter defines from _brand.yml ahead
% of this file; fonts are set by the filter from the brand's typography.
% Geometry: 16:9 frame of 16 cm x 9 cm. 1 cm here is 80 px in the live deck.

\usepackage{adjustbox}      % \lkfigureh: a maximum height that never upscales
\usepackage{tikz}
\usetikzlibrary{arrows.meta,positioning,calc,fit,backgrounds,decorations.pathreplacing}
\usepackage{etoolbox}
\usepackage{ragged2e}

% the brand's tokens, written by filters/tokens.lua before xelatex runs
\InputIfFileExists{_lecturekit/tokens.tex}{}{}

% fall-backs for a deck without a brand file
\providecolor{nmInk}{HTML}{202B33}
\providecolor{nmMuted}{HTML}{566573}
\providecolor{nmHairline}{HTML}{C8C3BA}
\providecolor{nmPaper}{HTML}{F3F2EE}
\providecolor{nmSlide}{HTML}{FFFFFF}
\providecolor{nmConcept}{HTML}{175B8E}
\providecolor{nmCaveat}{HTML}{A33A32}
\providecolor{nmExample}{HTML}{956000}
\providecolor{nmAnchor}{HTML}{006C65}
\providecolor{nmQuestion}{HTML}{70508C}
\providecolor{nmPlain}{HTML}{566573}
% optional tokens: heading (titles) and body (running text) fall back to ink
\providecolor{nmHeading}{named}{nmInk}
\providecolor{nmBody}{named}{nmInk}

\providecommand{\lkheadfont}{\sffamily}

% ---------------------------------------------------------------- colours
\setbeamercolor{normal text}{fg=nmBody,bg=nmSlide}
\setbeamercolor{background canvas}{bg=nmSlide}
\setbeamercolor{structure}{fg=nmInk}
\setbeamercolor{frametitle}{fg=nmHeading,bg=}
\setbeamercolor{title}{fg=nmHeading}
\setbeamercolor{subtitle}{fg=nmMuted}
\setbeamercolor{author}{fg=nmBody}
\setbeamercolor{institute}{fg=nmBody}
\setbeamercolor{date}{fg=nmMuted}
\setbeamercolor{itemize item}{fg=nmMuted}
\setbeamercolor{itemize subitem}{fg=nmMuted}
\setbeamercolor{enumerate item}{fg=nmMuted}
\setbeamercolor{caption name}{fg=nmMuted}
\setbeamercolor{footline}{fg=nmMuted}
\setbeamercolor{block title}{fg=nmInk,bg=}
\setbeamercolor{block body}{fg=nmBody,bg=}
% strong text in ink (the emphasis colour), as in the live deck
\renewcommand{\textbf}[1]{{\color{nmInk}\bfseries #1}}
\setbeamercolor{alerted text}{fg=nmCaveat}
\setbeamercolor{example text}{fg=nmExample}
\hypersetup{colorlinks=true,linkcolor=nmConcept,urlcolor=nmConcept,citecolor=nmConcept}

% ---------------------------------------------------------------- fonts
\setbeamerfont{frametitle}{size=\Large,series=\mdseries,family=\lkheadfont}
\setbeamerfont{title}{size=\LARGE,series=\mdseries,family=\lkheadfont}
\setbeamerfont{subtitle}{size=\normalsize,series=\mdseries}
\setbeamerfont{author}{size=\small}
\setbeamerfont{institute}{size=\small}
\setbeamerfont{date}{size=\small}
\setbeamerfont{footline}{size=\tiny}
\setbeamerfont{caption}{size=\scriptsize}
\setbeamerfont{block title}{series=\bfseries,size=\normalsize}

% ---------------------------------------------------------------- templates
\setbeamertemplate{navigation symbols}{}
\setbeamertemplate{itemize items}{\raisebox{0.18ex}{\scriptsize$\bullet$}}
\setbeamertemplate{itemize subitem}{\raisebox{0.2ex}{\tiny$\circ$}}
\setbeamertemplate{blocks}[default]
\setbeamertemplate{caption}{\insertcaption}
\setbeamertemplate{bibliography item}{}
\setbeamersize{text margin left=0.7cm,text margin right=0.7cm}

% the eyebrow: a role set by \lkrole{<colour>}{<label>} inside the frame;
% the frame title is typeset when the frame ends, so the body may set it
\newcommand{\lkCurrentRoleColor}{}
\newcommand{\lkCurrentRoleLabel}{}
\newcommand{\lkrole}[2]{%
  \gdef\lkCurrentRoleColor{#1}%
  \gdef\lkCurrentRoleLabel{#2}%
}
\BeforeBeginEnvironment{frame}{\gdef\lkCurrentRoleColor{}\gdef\lkCurrentRoleLabel{}}

\setbeamertemplate{frametitle}{%
  \vskip 0.45cm%
  \ifdefempty{\lkCurrentRoleLabel}{%
    \vphantom{\scriptsize X}\par\vskip 0.05cm%
  }{%
    {\scriptsize\bfseries\color{nm\lkCurrentRoleColor}%
      \addfontfeature{LetterSpace=12}\MakeUppercase{\lkCurrentRoleLabel}}\par\vskip 0.05cm%
  }%
  {\usebeamerfont{frametitle}\usebeamercolor[fg]{frametitle}\insertframetitle}\par%
  \vskip 0.15cm%
}
\setbeamertemplate{footline}{%
  \begin{beamercolorbox}[wd=\paperwidth,ht=0.45cm,dp=0.15cm,leftskip=0.7cm,rightskip=0.7cm]{footline}%
    {\color{nmHairline}\rule{\dimexpr\paperwidth-1.4cm\relax}{0.3pt}}\par\vskip 0.08cm%
    \usebeamerfont{footline}\insertshorttitle\hfill\insertframenumber\,/\,\inserttotalframenumber%
  \end{beamercolorbox}%
}

% ---------------------------------------------------------------- figures
% \lkfigureh{<max height>}{<file>}: as wide as the line allows, at most the
% given height, never upscaled (draw TikZ at display size).
\newcommand{\lkfigureh}[2]{%
  \begin{center}%
    \adjustbox{max height=#1,max width=\linewidth,center}{\includegraphics{#2}}%
  \end{center}%
}
\newcommand{\nmfigureh}[2]{\lkfigureh{#1}{#2}}
% a labelled placeholder so that a deck with a missing figure still builds
\newcommand{\lkplaceholder}[2]{%
  \begin{center}%
    \begin{tikzpicture}%
      \draw[nmCaveat,dashed,thin] (0,0) rectangle (0.8\linewidth,#1);%
      \node[nmCaveat,font=\scriptsize] at (0.4\linewidth,0.5*#1) {#2};%
    \end{tikzpicture}%
  \end{center}%
}

% role colour utility spans: [text]{.concept} from Pandoc arrive as \textcolor
% via Quarto; these macros exist for hand-written LaTeX in slides
\newcommand{\lkconcept}[1]{{\color{nmConcept}#1}}
\newcommand{\lkcaveat}[1]{{\color{nmCaveat}#1}}
\newcommand{\lkexample}[1]{{\color{nmExample}#1}}
\newcommand{\lkanchor}[1]{{\color{nmAnchor}#1}}
\newcommand{\lkquestion}[1]{{\color{nmQuestion}#1}}

% speaker notes stay out of the study PDF unless asked for
\setbeameroption{hide notes}

% tight lists, ragged right
\RaggedRight
\setlength{\parskip}{0.3em}
\setlength{\leftmargini}{1.2em}
